% =====================================================================
% Kurzfassung (German), Abstract (English), acknowledgements
% Rule of thumb: half a page to a page each, no sources, no figures.
% =====================================================================

\kurzfassung{%
Die Kurzfassung beschreibt in wenigen Absätzen, worum es in der Arbeit geht:
Ausgangsproblem, gewählter Lösungsansatz, umgesetztes Ergebnis und der Nutzen
für die Anwenderinnen und Anwender. Sie muss ohne den Rest der Arbeit
verständlich sein.

Der zweite Absatz beschreibt üblicherweise die technische Umsetzung auf
Überblicksebene: welche Komponenten entstanden sind und wie sie zusammenspielen.

Der letzte Absatz fasst das Ergebnis zusammen und benennt den erzielten
Mehrwert.

Die Spielanalyse im Fussball (Laufdistanzen, Positionsdaten und taktischen
Auswertungen) kostet oft Geld und viel Hardware. Kleine Vereine können sich diese Funktion oft nicht leisten. Das Diplomprojekt
VerTex setzt hier an. Es ist ein System das aus Videoaufnahmen des Spielfelds Bewegungs- und Spielstatistiken
gewinnt und diese Trainerinnen und Trainern ohne technisches Vorwissen
zugänglich macht.

Das Spielfeld wird mit einem synchronisierten Zwei-Kamera-Setup erfasst und zu
einem Panorama zusammengefügt. Ein Computer-Vision-Modell erkennt Spieler und Ball und verfolgt sie über die ganze Spiellänge. 
Mittels einer Homogprahie werden die Bildkoordinaten in Spielfeldkoordinaten umgerechnet. 
Die Bewegungsdaten werden über eine API in einer Datenbank gespeichert. 
Aus den Daten werden Laufwege, Distanzen und Teamstatistiken berechnet. 
Eine Web-Applikation greift über die API auf die Daten zu und stellt sie interaktiv dar.

Als Ergebnis soll eine interaktive Web-Applikation zur Verfügung gestellt werden,
die die Auswertung von Heatmaps und Teamstatistiken visuell darstellt. 
Weiters, soll ein fine-getuntes Computer-Vision-Modell für fußballspezifische Objekterkennung entwickelt werden. 




}

\abstracttext{%
The abstract is the English translation of the Kurzfassung. It covers the same
content in the same structure: the initial problem, the chosen approach, the
implemented result, and the benefit for its users.

The second paragraph outlines the technical realisation at an overview level:
which components were built and how they interact.

The final paragraph summarises the outcome and states the value it creates.

Match analysis in football (running distances, position data and tactical
evaluations) often requires a lot of money and hardware. Small clubs frequently
cannot afford this functionality. The diploma project VerTex addresses this
problem. It is a system that derives movement and match statistics from video
recordings of the pitch and makes them accessible to coaches without any
technical background.

The pitch is captured with a synchronised two-camera setup, and the images are
stitched into a panorama. A computer vision model detects players and the ball
and tracks them over the entire duration of the match. Using a homography, the
image coordinates are converted into pitch coordinates. The movement data is
stored in a database via an API. Running paths, distances and team statistics
are calculated from this data. A web application accesses the data through the
API and presents it interactively.

As a result, an interactive web application is to be provided that visualises
the evaluation of heatmaps and team statistics. Furthermore, a fine-tuned
computer vision model for football-specific object detection is to be
developed.
}

\danksagung{%
An dieser Stelle bedanken wir uns bei allen, die uns während der Erstellung
dieser Diplomarbeit unterstützt haben.

Unser besonderer Dank gilt unseren Betreuerinnen und Betreuern für die
fachliche Unterstützung und die investierte Zeit.
}
